\section{Corank modulo $p$}\label{sec:content}

This section proves the gain in the denominators. The principle is that of the arithmetic method of Hata
and Rhin--Viola, adapted to a determinant. If the reduction modulo $p$ of $p^2\cA^0_N$ has corank $e$, then
$\Delta_N(z)$ gains a factor $p^e$ at every $z$ with denominator prime to $p$
(Proposition~\ref{prop:valuation}). So we need explicit row combinations $c\in\F_p^n$ with
$c^{\mathsf T}\cA^0_{N,p}=0$.

Here is the heuristic. Modulo $p$ the Chebyshev polynomials of index $kp+d$ factor through $X^p$, so rows
placed symmetrically about the centres $g+kp$ are related. Suitable combinations of such rows have row
polynomials with long runs of vanishing base-$p$ digits. By the digit formula their values at most moment
columns then vanish, and the few remaining values are removed by imposing linear conditions. Each surviving
dimension is a unit of corank.

We first prove the counting inequality (Theorem~\ref{thm:count}). It says that a space of row combinations
subject to a few linear conditions lies in the left kernel. We then derive four closed formulas from it,
which give most of the gain. Finally we describe the further families and the resulting function $\kappa$.
Three facts about characteristic $p$ are used throughout: the Chebyshev identities
\begin{equation}\label{eq:cheb-p}
 \begin{gathered}
 T_p=X^p,\qquad U_{p-1}=(X^2-1)^{(p-1)/2},\\
 T_{kp}=T_k(X^p),\qquad U_{kp-1}=(X^2-1)^{(p-1)/2}\,U_{k-1}(X^p)
 \end{gathered}
\end{equation}
in $\F_p[X]$, the relation $(1-X)^p=1-X^p$, and the digit formula \eqref{eq:column-digits}. The last two
identities of \eqref{eq:cheb-p} follow from $T_{kp}=T_k\circ T_p$ and $U_{kp-1}=(U_{k-1}\circ T_p)\,U_{p-1}$.

\subsection{Setting}
Fix parameters, an integer $N\ge1$ and an odd prime $p$ with $H<p^2$. For $0\le r<n$ put $m_r=r-g$, so that
$-g\le m_r\le n-1-g$ and $n-1-g\ge g$. In the notation of \eqref{eq:rows},
\[
 P_r=(1-X)^h\operatorname{rev}_{C-1}T_{m_r},\qquad D_r=(1-X)^h\operatorname{rev}_{C-1}U_{m_r-1},
\]
where $T_{-m}=T_m$ and $U_{-m-1}=-U_{m-1}$, so that $D_g=0$. For $c\in\F_p^n$ write
$S_c=\sum_rc_r\operatorname{rev}_{C-1}T_{m_r}$ and $U_c=\sum_rc_r\operatorname{rev}_{C-1}U_{m_r-1}$, so that
$P_c=(1-X)^hS_c$ and $D_c=(1-X)^hU_c$. The lowest degree of $\operatorname{rev}_{C-1}T_m$ is at least
$C-1-|m|\ge A=b+q+2g$, and that of $\operatorname{rev}_{C-1}U_{m-1}$ is at least $A+1$. Put $\chi=g+q+h$. Then
the polynomials $Q_c$ and $\widehat D_c$ of \eqref{eq:Qc} have their coefficients in $[A+1,\,L+\chi+p-1]$
and $[A+1,\,L+\chi-1]$ respectively.

A vector $c$ lies in the left kernel of $\cA^0_{N,p}$ as soon as the right side of
\eqref{eq:column-digits} vanishes for every column $k$. We only use columns whose $q+1$ moment columns
share the same base-$p$ digit.

\subsection{The counting inequality}
All the kernels of this section arise in the same way. A space of row combinations is chosen, a few linear
conditions are imposed, and the conditions make every value vanish. The dimension of the space minus the
number of conditions is then a lower bound for the corank. The following theorem states this once. It also
allows $q\ge p$.

\begin{theorem}\label{thm:count}
Let $N\ge1$ and let $p$ be an odd prime with $H<p^2$. Write $q=ap+r$ with $0\le r<p$. Let $\cV$ be a
finite-dimensional vector space over $\F_p$, $\iota\colon\cV\to\F_p^n$ an injective linear map and $\Phi$ a
linear map on $\cV$. Suppose that for every $x\in\ker\Phi$ there is a polynomial $g_x\in\F_p[X]$ of degree
less than $r$ with
\[
 \sum_{v=0}^a(-1)^{a-v}\binom av\,R_{\iota(x)}(j+vp)=g_x(j)\qquad(b\le j<L-ap).
\]
Then $\nu_N(p)\ge\dim\cV-\operatorname{rank}\Phi$.
\end{theorem}

For $q<p$ the hypothesis says that the conditions $\Phi$ force all values at the moment columns to vanish.

\begin{proof}
By \eqref{eq:rational-part}, the entry of $c^{\mathsf T}\cA^0_{N,p}$ at column $k$ is
$\sum_{v=0}^q(-1)^v\binom qvR_c(b+k+v)$. This is the polynomial $(1-T)^q$ applied to $R_c$, where $T$ acts as
the shift. Modulo $p$ we have $(1-T)^p=1-T^p$, hence
\[
 (1-T)^q=(1-T)^r(1-T^p)^a .
\]
This is Lucas's congruence for the binomial coefficients $\binom qv$. So the entry at $k$ equals
$\sum_{u=0}^r(-1)^u\binom ru\,\Sigma_c(b+k+u)$, where
$\Sigma_c(j)=\sum_{v=0}^a(-1)^v\binom av R_c(j+vp)$. The arguments satisfy $b\le j$ and
$j+ap\le b+n-1+q=L-1$. For $c=\iota(x)$ with $x\in\ker\Phi$ we have $\Sigma_c(j)=\pm g_x(j)$. The difference
of order $r$ of a polynomial of degree less than $r$ vanishes identically. So $\iota(\ker\Phi)$ lies in the
left kernel of $\cA^0_{N,p}$. Since $\iota$ is injective, $\nu_N(p)\ge\dim\ker\Phi=\dim\cV-\operatorname{rank}\Phi$.
\end{proof}

Suppose instead that the conditions only force the values on a set $\mathcal C$ of columns to vanish. The
same argument, with the remaining $n-|\mathcal C|$ entries added to the conditions, gives
$\nu_N(p)\ge\dim\cV-\operatorname{rank}\Phi-(n-|\mathcal C|)$. The proofs of the four closed formulas below
are of this form. In Theorems~\ref{thm:T} and~\ref{thm:U} the set $\mathcal C$ consists of all columns, and in
Theorems~\ref{thm:U6} and~\ref{thm:S} it is a block of columns.

\subsection{Four closed formulas}
Each closed formula exhibits explicit vectors in the left kernel. The rows used are symmetric about the
centres $g+kp$, and the identities \eqref{eq:cheb-p} make their combinations vanish modulo $p$. Recall
$e=(p+1)/2-q-h$.

\begin{theorem}\label{thm:content}
Let $N\ge1$ and let $p$ be an odd prime with $H<p^2$. Then $\nu_N(p)$ is at least each of the following
quantities whenever the stated conditions hold. Negative values are replaced by $0$.
\begin{enumerate}[label=(\alph*)]
\item (Theorem~\ref{thm:T}) If $2g<p$, $H\le4p$ and $k_T=\min(g,\,n-1-g-2p,\,C-3p)\ge1$:
\[
 k_T-(n+q+g+h-3p)_+ .
\]
\item (Theorem~\ref{thm:U}) If $2h\le p+1$, $L\le3p\le H$, $n+q+g+h\le3p$ and $2p\le n+q$:
\[
 \min(n,\,2p+2g+1)-g-2p-1 .
\]
\item (Theorem~\ref{thm:U6}) If $e\ge0$, $n\le4p+2g+1$, $L\le6p$ and $b+q\le p+\varsigma_0$, with
$\varsigma_0=\max(q,H+q-6p)$:
\[
 \min\Bigl(\frac{p-1}2-g,\ n+b+q+h-\frac{11p+1}2\Bigr).
\]
\item (Theorem~\ref{thm:S}) If $e\ge0$ and $n\le4p+2g+1$, then for every integer $k\ge2$ the quantity $m_1-m_0-e$, with
\[
 m_0=\max(b+q,\,L+\chi-kp),\qquad m_1=\min(p,\,L-(k-1)p,\,A+1).
\]
\end{enumerate}
\end{theorem}

For our parameters these bounds apply for $\frac{74}5N<p\le\frac{954}{29}N$, where they give most of the gain
(Section~\ref{ssec:kappa}). We prove the four parts in turn.

\subsection{Antisymmetric combinations of rows}
\begin{theorem}\label{thm:T}
Let $2g<p$, $H\le4p$ and $k_T=\min(g,\,n-1-g-2p,\,C-3p)\ge1$. Then
$\nu_N(p)\ge k_T-(n+q+g+h-3p)_+$.
\end{theorem}

The rows used are the differences of rows placed symmetrically about the three centres $g$, $g+p$ and
$g+2p$.

\begin{proof}
For $1\le d\le k_T$ and $k=0,1,2$ let $\epsilon_k(d)=e_{g+kp+d}-e_{g+kp-d}\in\F_p^n$. These rows exist
because $d\le g$ and $g+2p+d\le n-1$. The $3k_T$ vectors $\epsilon_k(d)$ are linearly independent, since their
supports are disjoint. Indeed the supports lie in the intervals $[kp,kp+2g]$ of row indices, which are
disjoint because $2g<p$. The row $g+kp\pm d$ has $m=kp\pm d$. By $T_{a+b}-T_{a-b}=2(X^2-1)U_{a-1}U_{b-1}$
and $U_{a+b-1}-U_{a-b-1}=2T_aU_{b-1}$ together with \eqref{eq:cheb-p},
\begin{align*}
 T_{kp+d}-T_{kp-d}&=2(X^2-1)^{(p+1)/2}\,U_{k-1}(X^p)\,U_{d-1},\\
 U_{kp+d-1}-U_{kp-d-1}&=2\,T_k(X^p)\,U_{d-1}.
\end{align*}
Put $R_d=2(1-X)^hX^{C-d}\operatorname{rev}_{d-1}U_{d-1}$, the value of $D_{\epsilon_0(d)}$. Reverse at degree
$C-1$ and use $\operatorname{rev}_{kp}(T_k(X^p))=(\operatorname{rev}_kT_k)(X^p)$, with $T_0=1$, $T_1=X$ and
$T_2=2X^2-1$. One finds
\[
 D_{\epsilon_1(d)}=X^{-p}R_d,\qquad D_{\epsilon_2(d)}=2X^{-2p}R_d-R_d,
\]
\begin{gather*}
 XWP_{\epsilon_0(d)}=0,\qquad XWP_{\epsilon_1(d)}=X^{-p}R_d-X^pR_d,\\
 XWP_{\epsilon_2(d)}=2X^{-2p}R_d-2R_d .
\end{gather*}
The polynomial $R_d$ is supported in degrees $[C-d,H-1]$, which lies in $[3p,4p)$ because $d\le C-3p$ and
$H\le4p$. So at each residue $\ell$ it has a single base-$p$ digit $\varrho_\ell(d)$, in position $u=3$. For
$x=(x_0,x_1,x_2)\in\F_p^3$ the value of $\sum_kx_k\epsilon_k(d)$ at the moment column $Jp+\ell$ is therefore,
by Lemma~\ref{lem:digits}, $\varrho_\ell(d)\,f_J(x)$ with
\[
 f_J(x)=x_0\beta_{3,J}+x_1(\alpha_{2,J}-\alpha_{4,J}+\beta_{2,J})+x_2(2\alpha_{1,J}-2\alpha_{3,J}+2\beta_{1,J}-\beta_{3,J}).
\]
This depends neither on $d$ nor on $\ell$. The moment columns with $J\ge3$ do not meet the support, since
$3p+\ell\ge C-d\ge C-g=L$. Choose $x^*\ne0$ with $f_1(x^*)=f_2(x^*)=0$, which is two linear equations in three
unknowns. For $c=\sum_ds_d\sum_kx^*_k\epsilon_k(d)$ the values at the columns with $J=1,2$ vanish. The columns
with $J=0$ have $\ell\ge b$, and the value there is $f_0(x^*)\sum_ds_d\varrho_\ell(d)$, the coefficient of
$X^{3p+\ell}$ in $\sum_ds_dR_d$. These are at most $(H-3p-b)_+=(n+q+g+h-3p)_+$ linear conditions on
$s=(s_d)$. The map $s\mapsto c$ is injective, so the solutions form a space of dimension at least
$k_T-(n+q+g+h-3p)_+$ in the left kernel.
\end{proof}

\subsection{A reflection of Chebyshev polynomials}
The next bound rests on a reflection formula modulo $p$.

\begin{lemma}\label{lem:u3-reflect}
Let $p$ be an odd prime and $\varpi=(p+1)/2$. In $\F_p[X]$, for every integer $m$,
\[
 T_{p-m}=X^pT_m-(X^2-1)^{\varpi}\,U_{m-1}.
\]
\end{lemma}

\begin{proof}
Both sides satisfy the three-term recurrence in $m$, so it suffices to check $m=0$ and $m=1$. For $m=0$ it
is $T_p=X^p$, by \eqref{eq:cheb-p}. For $m=1$ use $T_{p+1}=XT_p+(X^2-1)U_{p-1}$ and $T_{p-1}=2XT_p-T_{p+1}$.
\end{proof}

\begin{theorem}\label{thm:U}
Let $2h\le p+1$, $L\le3p\le H$, $n+q+g+h\le3p$ and $2p\le n+q$, and put $n^*=\min(n,2p+2g+1)$. Then
$\nu_N(p)\ge n^*-g-2p-1$.
\end{theorem}

\begin{proof}
We restrict $c$ to the rows $r<n^*$, so that $m_r\le n^*-1-g\le2p+g$. Apply Lemma~\ref{lem:u3-reflect} with
$m=m_r$ and reflect at degree $p+C-1$. For every such row this gives
\[
 (1-X^2)^\varpi\operatorname{rev}_{C-1}U_{m-1}-X\operatorname{rev}_{C-1}T_m=-X^{p+1}\operatorname{rev}_{C-1}T_{p-m}.
\]
The right side is divisible by $X^{p+C-|p-m|}$, hence by $X^L$ since $|p-m|\le p+g=p+C-L$. Multiplying by
$(1-X)^h$ and summing, we get for every $c$ supported on these rows
\begin{equation}\label{eq:identityK}
 X^L\mid(1-X^2)^\varpi D_c-XP_c .
\end{equation}

Let $\cV$ be the space of such $c$ with $(1-X^2)^\varpi\mid P_c$. Since $(1-X)^h$ divides every $P_r$, this is
$\varpi-h$ conditions at $X=1$ and $\varpi$ conditions at $X=-1$, that is, at most $p+1-h$ linear
conditions. For $c\in\cV$ write $P_c=(1-X^2)^\varpi Y$ and $\widetilde Q=XY$. Since $(1-X^2)^p=1-X^{2p}$,
\[
 XWP_c=(1-X^{2p})\,\widetilde Q .
\]
The polynomial $P_c$ has no coefficients below $X^{A^*}$, $A^*=C-n^*+g$, and none from $X^H$ on. So
$\widetilde Q$ is supported in $[A^*+1,\,H-p-1]$. Dividing \eqref{eq:identityK} by the unit $(1-X^2)^\varpi$
modulo $X^L$ gives $D_c[k]=\widetilde Q[k]$ for $k<L$.

Now let $H-3p\le\ell<L-2p$. The inequality $L-2p\le A^*+1$ is equivalent to $n^*\le2p+2g+1$, and
$L-2p\le p$ since $L\le3p$. By the supports, $\widetilde Q[\ell]=0$, $\widetilde Q[2p+\ell]=0$ and
$\widetilde Q[3p+\ell]=0$, and $D_c[3p+\ell]=0$ because $3p+\ell\ge H$. So with
$\varkappa=\widetilde Q[p+\ell]$ the base-$p$ digits of $XWP_c$ at the residue $\ell$ are
$(0,\varkappa,0,-\varkappa,0,\dots)$, and those of $D_c$ are $(0,\varkappa,0,0,\dots)$. Hence, by
Lemma~\ref{lem:digits}, the value at every moment column $Jp+\ell$ with such an $\ell$ is $\varkappa$ times a
constant depending on $J$ only. On $\cV$ these columns, which have $J\in\{0,1,2\}$, therefore span a space
of dimension at most one for each residue $\ell$.

Put $\ell_0=H-3p$ and $\ell_1=L-2p$. The hypotheses give $\ell_0\le b\le\ell_1$. For $\ell_0\le\ell<b$ the
moment columns $j\in[b,L)$ with this residue have $J\in\{1,2\}$, and for $b\le\ell<\ell_1$ they have
$J\in\{0,1,2\}$. So the saving is $2\ell_1-\ell_0-b$. The number of linear conditions needed to make all
values on the $n+q=L-b$ moment columns vanish on $\cV$ is therefore
$(p+1-h)+(n+q)-(2\ell_1-\ell_0-b)$. Substituting $\ell_1=L-2p$ and $\ell_0=H-3p$, this is $2p+1+g$. So the
space of solutions in the span of $n^*$ rows has dimension at least $n^*-g-2p-1$.
\end{proof}

\subsection{A block of columns: the digits $1$ to $5$}
\begin{theorem}\label{thm:U6}
Let $e=(p+1)/2-q-h\ge0$, $n\le4p+2g+1$ and $L\le6p$, and put $\varsigma_0=\max(q,H+q-6p)$ and
$\varsigma_1=L-5p$. Assume $b+q\le p+\varsigma_0$. Then $\nu_N(p)\ge\varsigma_1-\varsigma_0-e$, that is,
\[
 \nu_N(p)\ge\min\Bigl(\frac{p-1}2-g,\ n+b+q+h-\frac{11p+1}2\Bigr).
\]
The first term is the minimum if $6p\le H$ and the second if $6p\ge H$.
\end{theorem}

In Theorems~\ref{thm:T} and~\ref{thm:U} the conditions are imposed on the rows. Here the argument concerns a
block of columns. We first need an identity.

\begin{lemma}\label{lem:K6}
In $\F_p[X]$, for every integer $k$, $2X^p(X^2-1)^{(p-1)/2}T_k-(2X^{2p}-1)U_{k-1}=U_{2p-k-1}$. If
$n\le4p+2g+1$, then for every $c$, $X^L\mid 2XWS_c-(2-X^{2p})U_c$, and consequently
\begin{equation}\label{eq:K6}
 2\,Q_c[i]=2\,\widehat D_c[i]-\widehat D_c[i-2p]\qquad(i<L).
\end{equation}
\end{lemma}

\begin{proof}
For all integers $m,k$ one has $U_{m-1}T_k-T_mU_{k-1}=U_{m-k-1}$, since both sides satisfy the same
recurrence in $k$ and agree for $k=0,1$. Take $m=2p$ and use \eqref{eq:cheb-p}: $T_{2p}=2X^{2p}-1$ and
$U_{2p-1}=2T_pU_{p-1}=2X^p(X^2-1)^{(p-1)/2}$. Reflect at degree $2p+C-1$. The first term becomes
$2XW\operatorname{rev}_{C-1}T_k$, the second $(2-X^{2p})\operatorname{rev}_{C-1}U_{k-1}$, and the right side
$X^{2p+C-1}U_{2p-k-1}(1/X)$. For $k=m_r\ge-g$ with $2p-k-1\ge0$ the right side is divisible by $X^{C+k}$, hence
by $X^L$. If $2p-k-1<0$, it has degree $k-2p-1$ and is divisible by $X^{4p+C-k}$, hence by $X^L$ when
$k\le4p+g$. This holds because $k\le n-1-g\le4p+g$. Summing with weights $c_r$ and multiplying by
$(X-1)^q(1-X)^h$ gives \eqref{eq:K6}.
\end{proof}

\begin{proof}[Proof of Theorem~\ref{thm:U6}]
Let $\mathcal M=[\varsigma_0,\varsigma_1)$ and $w=\varsigma_1-\varsigma_0$. We may assume $w>e$. The block consists of the columns
$k(J,\varsigma)=Jp+\varsigma-(b+q)$ for $1\le J\le5$ and $\varsigma\in \mathcal M$. These indices are nonnegative by the hypothesis
$b+q\le p+\varsigma_0$, and smaller than $n$ because $5p+\varsigma_1-1=L-1$. There are $5w$ of them. As $\varsigma<p$, formula
\eqref{eq:column-digits} applies, and the value of $c$ at $k(J,\varsigma)$ is
$\sum_u\alpha_{u,J}Q_c[up+\varsigma]+\beta_{u,J}\widehat D_c[up+\varsigma]$.

\emph{Supports.} For $\varsigma\in \mathcal M$ we have $p+\varsigma\le p+\varsigma_1-1=L-4p-1\le A$ by $n\le4p+2g+1$. So the digits $u=0,1$
of $Q_c$ and $\widehat D_c$ vanish. From $\varsigma_0\ge H+q-6p=L+\chi-6p$ we get that the digits $u\ge6$ of
$\widehat D_c$ and $u\ge7$ of $Q_c$ vanish. For $u\le5$ we have $up+\varsigma\le L-1$, so \eqref{eq:K6} applies.

\emph{From $Q$ to $\widehat D$.} For $u=2,\dots,5$, \eqref{eq:K6} gives
$\widehat D_c[up+\varsigma]=Q_c[up+\varsigma]+\tfrac12\widehat D_c[(u-2)p+\varsigma]$, and for $u=2,3$ the last index is at most $A$,
so that term vanishes. Hence if $Q_c[up+\varsigma]=0$ for $u=2,\dots,5$ and all $\varsigma\in \mathcal M$, then by induction
$\widehat D_c[up+\varsigma]=0$ for $u=2,\dots,5$, and all digits of $\widehat D_c$ in $\mathcal M$ vanish.

\emph{A relation from the Frobenius map.} We have
$Q_c=(-1)^qX(1-X)^{q+h+(p-1)/2}(1+X)^{(p-1)/2}S_c$, and $q+h+(p-1)/2+e=p$. Since $(1-X)^p=1-X^p$,
\[
 (1-X)^e\,Q_c=(1-X^p)\,Y_c,\qquad Y_c=(-1)^qX(1+X)^{(p-1)/2}S_c,
\]
and $Y_c$ is supported in $[A+1,\,C+(p-1)/2]$. Compare the coefficients at $up+\varsigma$, subtract the
coefficient at $(u-1)p+\varsigma$, and sum over $u=2,\dots,6$. The right side telescopes to
$Y_c[6p+\varsigma]-Y_c[p+\varsigma]$, which vanishes for $\varsigma\ge\varsigma_0+e$. The second term vanishes because $p+\varsigma\le A$,
and the first because $6p+\varsigma\ge H+q+e=C+(p+1)/2>\deg Y_c$. So, if $Q_c[up+\varsigma']=0$ for $u=2,\dots,5$ and all
$\varsigma'\in \mathcal M$, then
\[
 \sum_{t=0}^e(-1)^t\binom et\,Q_c[6p+\varsigma-t]=0\qquad(\varsigma_0+e\le\varsigma<\varsigma_1).
\]
This is a recurrence with unit leading coefficient. So if $Q_c[6p+\varsigma]=0$ for the first $e$ values
$\varsigma=\varsigma_0,\dots,\varsigma_0+e-1$, then it vanishes for all $\varsigma\in \mathcal M$.

\emph{Counting.} Impose on $c$ that the values at the $n-5w$ columns outside the block vanish, that
$Q_c[up+\varsigma]=0$ for $2\le u\le5$ and $\varsigma\in \mathcal M$ ($4w$ conditions), and that $Q_c[6p+\varsigma]=0$ for the first $e$
values of $\varsigma$. By the three steps above all the values at the block columns vanish. These are at most
$n-5w+4w+e=n-(w-e)$ linear conditions, so the left kernel has dimension at least $w-e$. Finally, $w-e$
equals $\varsigma_1-q-e=n+b+q+h-(11p+1)/2$ if $\varsigma_0=q$, and $(p-1)/2-g$ if $\varsigma_0=H+q-6p$.
\end{proof}

\subsection{The digits $1$ to $k$}\label{ssec:S}
The proof of Theorem~\ref{thm:U6} uses five digits. A shorter block uses fewer. The count is the same for
every number of digits $k\ge2$, and we prove it once. The cases $k=4$ and $k=3$ both occur for our
parameters. The case $k=3$ gives the bounds on $(\frac{186}7N,28N]$.

\begin{theorem}\label{thm:S}
Let $e=(p+1)/2-q-h\ge0$ and $n\le4p+2g+1$. Let $k\ge2$ be an integer and put
\[
 m_0=\max(b+q,\,L+\chi-kp),\qquad m_1=\min\bigl(p,\,L-(k-1)p,\,A+1\bigr).
\]
Then $\nu_N(p)\ge m_1-m_0-e$.
\end{theorem}

The proof rests on a statement about digits only. Write $\mathcal M=[m_0,m_1)$. For $\varsigma\in\mathcal M$ and
$u\ge0$ put $Q_u=Q_c[up+\varsigma]$ and $D_u=\widehat D_c[up+\varsigma]$, with the polynomials of
\eqref{eq:Qc}. The relation \eqref{eq:K6} reads $2Q_u=2D_u-D_{u-2}$ for $up+\varsigma<L$, with $D_{u-2}=0$
for $u<2$. The Frobenius relation $(1-X)^eQ_c=(1-X^p)Y_c$ of the proof of Theorem~\ref{thm:U6} gives
\begin{equation}\label{eq:frob-digits}
 \sum_{t=0}^e(-1)^t\binom et\,Q_c[i-t]=Y_c[i]-Y_c[i-p]\qquad(i\ge0).
\end{equation}

\begin{lemma}\label{lem:digit-vanishing}
Let $w=m_1-m_0>e$ and assume the hypotheses of Theorem~\ref{thm:S}. Suppose that $D_j=0$ for
$1\le j<k$ and all $\varsigma\in\mathcal M$, and that $Q_k=0$ for the first $e$ residues
$\varsigma=m_0,\dots,m_0+e-1$. Then $Q_u=D_u=0$ for all $u\ge0$ and all $\varsigma\in\mathcal M$.
\end{lemma}

\begin{proof}
Since $m_1\le A+1$, we have $\varsigma\le A$ and the digits $u=0$ of $Q_c$ and $\widehat D_c$ vanish. Since
$\varsigma\ge L+\chi-kp$, the digits $u\ge k$ of $\widehat D_c$ and $u\ge k+1$ of $Q_c$ vanish. Indeed these
polynomials have no coefficients beyond $L+\chi-1$ and $L+\chi+p-1$. The digits $D_j$ with $1\le j<k$ vanish by
assumption. For $1\le j<k$ we have $jp+\varsigma\le(k-1)p+m_1-1<L$. So \eqref{eq:K6} applies. It gives
$2Q_j=2D_j-D_{j-2}=0$, because $D_{j-2}$ is either absent or a digit already known to vanish. Only $Q_k$
remains.

Let $\varsigma\ge m_0+e$. Apply \eqref{eq:frob-digits} at $i=up+\varsigma$ for $u=0,\dots,k+1$ and add. The right sides
telescope to $Y_c[(k+1)p+\varsigma]$. This coefficient vanishes. Indeed $Y_c$ has degree at most
$C+(p-1)/2$, while $(k+1)p+\varsigma\ge p+L+\chi>C+(p-1)/2$, because $p+q+h>(p-1)/2$. On the left, every residue
$\varsigma-t$ with $0\le t\le e$ lies in $\mathcal M$. So all digits vanish except the digit $k$, and
\[
 \sum_{t=0}^e(-1)^t\binom et\,Q_c[kp+\varsigma-t]=0 .
\]
The term $t=0$ has coefficient $1$. The first $e$ values of $Q_k$ vanish by assumption. Induction on $\varsigma$
shows $Q_k=0$ on all of $\mathcal M$.
\end{proof}

\begin{proof}[Proof of Theorem~\ref{thm:S}]
We may assume $w=m_1-m_0>e$. The columns $Jp+\varsigma-(b+q)$ with $0\le J<k$ and $\varsigma\in\mathcal M$ exist.
They are nonnegative since $\varsigma\ge b+q$, and smaller than $n$ since $(k-1)p+\varsigma<L$. They are
distinct since $\varsigma<p$. So there are $kw$ of them. By \eqref{eq:column-digits} the value of $c$ at such a
column is a combination of the digits $Q_u,D_u$ at the residue $\varsigma$. Impose on $c$ that the values at the
other $n-kw$ columns vanish, that $D_j=0$ for $1\le j<k$ and all $\varsigma\in\mathcal M$ ($(k-1)w$ conditions),
and that $Q_k=0$ for the first $e$ residues. By Lemma~\ref{lem:digit-vanishing} all digits at the residues of
$\mathcal M$ vanish. So the values at the $kw$ columns vanish as well. There are at most
$(n-kw)+(k-1)w+e=n-(w-e)$ conditions, and the left kernel has dimension at least $w-e$.
\end{proof}

The lemma is the digit-only core of the argument. The proof uses no determinant and no lower bound for $N$.


\subsection{Decompositions modulo $p$}
The further families organize the row combinations as follows.

\begin{definition}\label{def:slot}
A \emph{decomposition modulo $p$} of a family of row combinations consists of
\begin{enumerate}[label=(\alph*)]
\item a linear map $x\mapsto c(x)$ from $(\F_p^K)^t$ to $\F_p^n$, injective on a subspace of admissible
 $x=(x_1,\dots,x_t)$;
\item polynomials $\pi_1,\dots,\pi_t\in\F_p[X]$ of degree less than $e_{\max}$;
\item linear maps $\mathcal T,\mathcal G\colon\F_p^K\to\F_p[Y]$,
\end{enumerate}
such that for admissible $x$
\begin{gather*}
 (X-1)^hP_{c(x)}=\sum_d\pi_d(X)\,\mathcal T(x_d)(X^p),\\
 (X-1)^h(X^2-1)^{\frac{p-1}2}D_{c(x)}=\sum_d\pi_d(X)\,\mathcal G(x_d)(X^p).
\end{gather*}
\end{definition}

The identities \eqref{eq:cheb-p} and $(1-X)^p=1-X^p$ provide such decompositions. By the digit formula, the
value of $c(x)$ at the moment column $Jp+\ell$ is a linear form in the $x_d$:
\[
 \sum_{0\le e'<e_{\max},\ p\mid H-1-\ell-e'}\ \Theta_{J,\,(H-1-\ell-e')/p}\Bigl(\sum_d\pi_d[e']\,x_d\Bigr).
\]
Here $\pi_d[e']$ is the coefficient of $X^{e'}$, and each $\Theta_{J,s}$ is a linear form in $J$ and the digit
$s$, determined by $\mathcal T$ and $\mathcal G$.

\subsection{Further families}\label{ssec:families}
The following lower bounds for $\nu_N(p)$ are proved for all parameters, by explicit choices of the space,
the conditions and the columns in Theorem~\ref{thm:count} or its variant. Each is a piecewise linear function
of $p$ and the parameters, positive on intervals described by linear inequalities. We describe them in
words. The statements are those of the formalization.
\begin{enumerate}[label=(F\arabic*),wide]
\item \emph{Frobenius products.} The row combinations have polynomials that are the product of a fixed
factor, a free polynomial of bounded degree $d_P$ in the Chebyshev variable, and a polynomial $\Xi(X^p)$ with
exponents in an interval of $K$ consecutive values. The conditions require $\Xi$ to annihilate the
functionals seen by the columns. If $\#\text{seen}$ of these functionals are independent, the bound is
$(d_P+1)(K-\#\text{seen})$. For our parameters this gives $66N-2p-1$ on $(\frac{97}3N,\frac{954}{29}N]$.
\item \emph{A truncated form of Theorem~\ref{thm:U6}.} Without the hypothesis $n\le4p+2g+1$, the bound
$\varsigma_1-\varsigma_0-e_+-(n-4p-2g-1)_+$ holds whenever $L\le6p$ and $b+q\le p+\varsigma_0$. For our parameters it gives
$5p-74N$ on $(\frac{74}5N,15N]$ and $p-14N$ on $(15N,16N]$, up to a bounded term.
\item \emph{A tail of the antisymmetric family.} With the four centres $g+kp$ ($0\le k\le3$) and
$k_T'=n-1-g-3p$, the bound is $p-1-(2g+q+h)$, under $4p\le L$, $b\le p$, $k_T'\le g$, $k_T'\le(p-1)/2$,
$3p+1+g\le n$ and $4p+b\le H$. For our parameters this is $p-20N-1$ on $(\frac{222}{11}N,\frac{146}7N]$.
\item \emph{Two families of decompositions and their sum.} The first fixes a digit $s\le H/p$ and a range of
blocks of rows, with polynomials $\pi_d$ divisible by a power of $X-1$. The second has polynomials
$(X^p-1)\pi$ with $\pi$ supported in an interval of degrees. Under a separation condition their kernels are
independent when both apply, at the same digit $s$ or at the digits $s$ and $s-1$. Then the sum of the two
bounds is a lower bound. This is the one place where two bounds are added. For our parameters it gives
bounds on six intervals between $8.4N$ and $10.9N$.
\item \emph{Residue intervals and rational cells.} For $p\le q$ one chooses the row combinations residue by
residue. On an interval $[\alpha,\beta)$ of residues of constant combinatorial type the bound is
$\beta-\alpha+p-q$ (Theorem~L$^*$). For $p\ge h$ and for $p\le h$ one uses cells defined over $\Q$. A cell
fixes a window of the digits, a range of blocks of rows and an integer profile, such that the value
of the combination at the columns of the cell has a fixed difference over $\Q$. For a prime $p$ prime to the modulus of the cell
the bound is $(\mathrm{hi}-\mathrm{lo})-(p-q)$ in the first case (Theorem~C) and
$(\mathrm{hi}-\mathrm{lo})-(h-p)-(p-q)$ in the second (Theorem~C$'$), where $[\mathrm{lo},\mathrm{hi})$ is an interval of positions
compatible with the cell. For our parameters each of these theorems supplies selected subintervals, with gaps,
of an interval. Theorem~C$'$ gives bounds on $(\frac{20}7N,\frac{66}{23}N]$ and $(\frac{80}{27}N,3N]$.
Theorem~L$^*$ gives bounds on $(\frac{74}{21}N,\frac{18}5N]$, $(\frac{37}{10}N,\frac{72}{19}N]$ and
$(\frac{74}{19}N,4N]$. Theorem~C gives bounds on $(4N,\frac{37}9N]$, $(\frac{72}{17}N,\frac{74}{17}N]$,
$(\frac92N,\frac{87}{19}N]$ and $(\frac{97}{21}N,\frac{37}8N]$. On $(\frac{87}{19}N,\frac{97}{21}N]$ the bound
combines two cells of Theorem~C and has the value $\mathrm{hi}_A+(\mathrm{hi}_B-\mathrm{lo}_B)-(p-4N)$, with
$\mathrm{lo}_B=97N-21p$, $\mathrm{hi}_A=\min(97N-21p,10N-2p)$ and $\mathrm{hi}_B=\min(10N-2p,70N-1-15p)$.
Intervals shifted by tiny rational amounts are listed in the formalization.
\item \emph{Affine cells, Theorem~HY$^+$ and Theorem~MT.} On $(\frac{83}{31}N,\frac{101}{37}N]$ and
$(\frac{36}{13}N,\frac{14}5N]$ the same
counting applies with cells on which the values of the combinations at consecutive digits differ by a fixed
slope, for primes $p>44$ prime to a fixed integer. Between these two intervals the bounds come from explicit
subspaces. Theorem~HY$^+$ uses the residue sums for the digits $s=11,10,9$ on the following intervals. For $s=11$
it gives $(\frac{70}9N,8N]$. For $s=10$ it gives $(\frac{17}2N,\frac{97}{11}N]$, and a variant with a short fold gives
$(\frac{97}{11}N,\frac{4409091}{500000}N]$. For $s=9$ it gives $(\frac{66}7N,\frac{202}{21}N]$, and a variant with a short
fold gives $(\frac{97}{10}N,\frac{8300000001}{850000000}N]$. A mixed family of antisymmetric and symmetric pairs
of rows (Theorem~MT) gives $\nu_N(p)\ge\frac{29}2p-77N-2$ on $(\frac{43}8N,\frac{97}{18}N]$.
\item \emph{Theorem Z.} Two further families, called the laws $\mathrm{L}$ and $\mathrm{H}$ in the formalization,
give bounds on $(\frac{186}{11}N,\frac{160}9N]$ and on $(\frac{132}7N,\frac{222}{11}N]$. Write $T=H-5p$ and $h_p=(p-1)/2$.
The law $\mathrm L$ gives
\begin{multline*}
 2\nu_N(p)\ge11p-2(n+b+q+g)-1\\
 {}+2\min(0,L-5p)+2\min(0,T-b)+2\min(0,h+h_p-T).
\end{multline*}
In the regime $5p\le L$ and $b\le T\le h+h_p$ the three minima vanish. For our parameters the three selected lines
are $\frac{11}2p-93N$, $\frac p2-7N$ and $80N-\frac92p$, up to a loss of at most two. The two laws count the conditions
that the Taylor coefficients of the combinations impose on a space of rows.
\end{enumerate}
The interval $(\frac{96666667}{10000000}N,\frac{290999999}{30000000}N]$ belongs to an explicit subspace of
Appendix~\ref{app:windows}, with $\kappa(\theta)=\frac{599999999}{100000000}-\frac\theta2$. It is not assigned to
Theorem~HY$^+$. The interval $(\frac{202}{21}N,\frac{29}3N]$ before it belongs to (F4).

In the cases that need a prime to avoid finitely many bad values, only $p$ larger than fixed integers occur,
and the corresponding $N$ are large.

For the primes below $12N$ that none of these closed forms covers, we use the counting inequality in a special
form. The space of row combinations is the tensor product of a space of polynomials with a fixed rational
subspace $V\subset\Q^s$, and the conditions are Taylor coefficients composed with fixed linear maps. The
subspaces $V$ are explicit but have no closed form. Appendix~\ref{app:windows} states this form
(Theorem~\ref{thm:profile}) and describes how its hypotheses are verified.

\subsection{The corank profile for our parameters}\label{ssec:kappa}
The bounds combine into the following statement.

\begin{theorem}\label{thm:profile-D}
For the parameters $(72,11,4,6,4)$ there are $148$ intervals $(a_i,b_i]$ with rational endpoints,
$\frac{185000001}{74000000}=a_1<b_1\le a_2<\dots<b_{148}=\frac{954}{29}$, and a function $\kappa$, affine on each
interval and $0$ outside their union, with the following property. For all large $N$ and every prime $p$ with
$a_iN<p\le b_iN$ for some $i$,
\[
 \nu_N(p)\ge N\kappa(p/N)-2 .
\]
On each interval the bound is given by one of the families of this section, or by the explicit subspaces of
Appendix~\ref{app:windows}. The area of $\kappa$ is the rational number
\begin{multline*}
 \int_0^\infty\kappa(\theta)\,d\theta\\
 =\frac{8933203314101219901240655133431309756867}{258903586359571936582800000000000000000}=\Chi\ldots .
\end{multline*}
\end{theorem}

The first endpoint exceeds $\frac52$ by $\frac{1}{74\cdot10^6}$. Some other endpoints and intercepts are moved by amounts below $10^{-7}$,
to absorb the integer rounding of the families. The intervals are listed in the formalization.

The contributions by ranges of $\theta=p/N$ are as follows. The third column counts the intervals on which
the bound is that of Theorem~\ref{thm:profile} with an explicit subspace.
\begin{center}
\begin{tabular}{@{}lrrl@{}}
\toprule
range of $\theta$ & intervals & with explicit subspace & area \\
\midrule
$(5/2,\,3]$ & $29$ & $9$ & $0.03074$ \\
$(3,\,6]$ & $67$ & $20$ & $0.36842$ \\
$(6,\,12]$ & $31$ & $13$ & $1.06902$ \\
$(12,\,954/29]$ & $21$ & $0$ & $33.03580$ \\
\midrule
total & $148$ & $42$ & $\Chi$ \\
\bottomrule
\end{tabular}
\end{center}
Almost all of the area lies in $(12,\frac{954}{29}]$. There the intervals with $\theta>\frac{74}5$ are as follows.
Adjacent intervals with the same line are merged.
\begin{center}
\begin{tabular}{@{}lll@{}}
\toprule
interval of $\theta$ & $\kappa(\theta)$ & bound \\
\midrule
$(74/5,\,15]$ & $5\theta-74$ & (F2) \\
$(15,\,16]$ & $\theta-14$ & (F2) \\
$(16,\,97/6]$ & $\theta/2-6$ & Theorem~\ref{thm:content}(c), $(p-1)/2-g$ \\
$(97/6,\,182/11]$ & $91-11\theta/2$ & Theorem~\ref{thm:content}(c), second term \\
$(186/11,\,86/5]$ & $11\theta/2-93$ & (F7), law $\mathrm L$ \\
$(86/5,\,87/5]$ & $\theta/2-7$ & (F7), law $\mathrm L$ \\
$(87/5,\,160/9]$ & $80-9\theta/2$ & (F7), law $\mathrm L$ \\
$(132/7,\,97/5]$ & $7\theta/2-66$ & (F7), law $\mathrm H$ with cap \\
$(97/5,\,20]$ & $31-3\theta/2$ & (F7), law $\mathrm H$ \\
$(20,\,222/11]$ & $91-9\theta/2$ & (F7), law $\mathrm H$ \\
$(222/11,\,146/7]$ & $\theta-20$ & (F3) \\
$(146/7,\,43/2]$ & $9\theta/2-93$ & Theorem~\ref{thm:content}(d), $k=4$ \\
$(43/2,\,87/4]$ & $\theta/2-7$ & Theorem~\ref{thm:content}(d), $k=4$ \\
$(87/4,\,160/7]$ & $80-7\theta/2$ & Theorem~\ref{thm:content}(d), $k=4$ \\
$(186/7,\,27]$ & $7\theta/2-93$ & Theorem~\ref{thm:content}(d), $k=3$ \\
$(27,\,28]$ & $5\theta/2-66$ & Theorem~\ref{thm:content}(d), $k=3$ \\
$(28,\,86/3]$ & $3\theta-80$ & Theorem~\ref{thm:content}(a), $3p-(n+q+h)$ \\
$(86/3,\,30]$ & $6$ & Theorem~\ref{thm:content}(a) up to $\theta=29$, then (b) \\
$(30,\,97/3]$ & $66-2\theta$ & Theorem~\ref{thm:content}(b), $n-g-2p$ \\
$(97/3,\,954/29]$ & $66-2\theta$ & (F1) \\
\bottomrule
\end{tabular}
\end{center}
The areas of these rows add up to $33.03580\ldots$, and the $127$ intervals below $12$
contribute $1.46818\ldots$. The exact sum is computed in the
formalization from the list of intervals. Figure~\ref{fig:corank} shows $\kappa$ on $(\frac{74}5,\frac{954}{29}]$.

\begin{figure}[ht]
\centering
\setlength{\unitlength}{1cm}
\begin{picture}(12.4,5.4)(-0.8,-0.9)
\put(0,0){\line(1,0){11.8}}
\put(0,0){\line(0,1){4.9}}
\put(0.000,0){\line(0,-1){0.12}}
\put(0.000,-0.4){\makebox(0,0){$14$}}
\put(1.200,0){\line(0,-1){0.12}}
\put(1.200,-0.4){\makebox(0,0){$16$}}
\put(2.400,0){\line(0,-1){0.12}}
\put(2.400,-0.4){\makebox(0,0){$18$}}
\put(3.600,0){\line(0,-1){0.12}}
\put(3.600,-0.4){\makebox(0,0){$20$}}
\put(4.800,0){\line(0,-1){0.12}}
\put(4.800,-0.4){\makebox(0,0){$22$}}
\put(6.000,0){\line(0,-1){0.12}}
\put(6.000,-0.4){\makebox(0,0){$24$}}
\put(7.200,0){\line(0,-1){0.12}}
\put(7.200,-0.4){\makebox(0,0){$26$}}
\put(8.400,0){\line(0,-1){0.12}}
\put(8.400,-0.4){\makebox(0,0){$28$}}
\put(9.600,0){\line(0,-1){0.12}}
\put(9.600,-0.4){\makebox(0,0){$30$}}
\put(10.800,0){\line(0,-1){0.12}}
\put(10.800,-0.4){\makebox(0,0){$32$}}
\put(0,0.000){\line(-1,0){0.12}}
\put(-0.35,0.000){\makebox(0,0)[r]{$0$}}
\put(0,0.750){\line(-1,0){0.12}}
\put(-0.35,0.750){\makebox(0,0)[r]{$1$}}
\put(0,1.500){\line(-1,0){0.12}}
\put(-0.35,1.500){\makebox(0,0)[r]{$2$}}
\put(0,2.250){\line(-1,0){0.12}}
\put(-0.35,2.250){\makebox(0,0)[r]{$3$}}
\put(0,3.000){\line(-1,0){0.12}}
\put(-0.35,3.000){\makebox(0,0)[r]{$4$}}
\put(0,3.750){\line(-1,0){0.12}}
\put(-0.35,3.750){\makebox(0,0)[r]{$5$}}
\put(0,4.500){\line(-1,0){0.12}}
\put(-0.35,4.500){\makebox(0,0)[r]{$6$}}
\qbezier(0.480,0.000)(0.540,0.375)(0.600,0.750)
\qbezier(0.600,0.750)(0.900,1.125)(1.200,1.500)
\qbezier(1.200,1.500)(1.250,1.531)(1.300,1.562)
\qbezier(1.300,1.562)(1.414,0.781)(1.527,0.000)
\qbezier(1.745,0.000)(1.833,0.600)(1.920,1.200)
\qbezier(1.920,1.200)(1.980,1.237)(2.040,1.275)
\qbezier(2.040,1.275)(2.153,0.637)(2.267,0.000)
\qbezier(2.914,0.000)(3.077,0.712)(3.240,1.425)
\qbezier(3.240,1.425)(3.420,1.087)(3.600,0.750)
\qbezier(3.600,0.750)(3.655,0.443)(3.709,0.136)
\qbezier(3.709,0.136)(3.912,0.390)(4.114,0.643)
\qbezier(4.114,0.643)(4.307,1.728)(4.500,2.812)
\qbezier(4.500,2.812)(4.575,2.859)(4.650,2.906)
\qbezier(4.650,2.906)(4.982,1.453)(5.314,0.000)
\qbezier(7.543,0.000)(7.671,0.562)(7.800,1.125)
\qbezier(7.800,1.125)(8.100,2.062)(8.400,3.000)
\qbezier(8.400,3.000)(8.600,3.750)(8.800,4.500)
\qbezier(8.800,4.500)(8.900,4.500)(9.000,4.500)
\qbezier(9.000,4.500)(9.300,4.500)(9.600,4.500)
\qbezier(9.600,4.500)(10.300,2.750)(11.000,1.000)
\qbezier(11.000,1.000)(11.169,0.578)(11.338,0.155)
\put(11.95,-0.4){\makebox(0,0)[l]{$\theta$}}
\end{picture}

\caption{The corank profile $\kappa(\theta)$ on $(\frac{74}5,\frac{954}{29}]$, where $\nu_N(p)\ge N\kappa(p/N)-2$ for
$p=\theta N$. The bound on the pole count is $f(\theta)=144$ on this whole range, so $\kappa$ is the gain over
OpenAI's denominators.}\label{fig:corank}
\end{figure}

On each interval we use a single bound. The families come from different constructions, and we have not
shown that their kernels are independent, except for the pair (F4). So the bounds are not added at a prime
where several apply.

For Theorem~\ref{thm:content} one checks the hypotheses on the intervals above. The hypothesis $e\ge0$
reads $\theta\ge16$, the hypothesis $n\le4p+2g+1$ reads $\theta\ge15$, and $L\le6p$ reads $\theta\ge29/2$. For
$16\le\theta\le97/6$ one has $\varsigma_0=H+q-6p=101N-6p$, so the first branch of Theorem~\ref{thm:U6}
applies. For larger $\theta$ one has $\varsigma_0=4N$ and the second branch applies, which is positive until
$\theta=182/11$. In Theorem~\ref{thm:S} with $k=4$, $m_1=p$ for $\theta\le43/2$ and $m_1=L-3p=87N-3p$ above, and
$m_0=\max(15N,101N-4p)$, so that $m_0=101N-4p$ for $\theta\le43/2$. The three values $m_1-m_0-e$ are
$9\theta/2-93$, $\theta/2-7$ and $80-7\theta/2$. With $k=3$ one has $m_0=101N-3p$ and
$m_1=\min(p,27N+1)$, and the value $m_1-m_0-e$ is $7\theta/2-93$ on $(\frac{186}7,27]$ and $5\theta/2-66$ on
$(27,28]$, up to the constants $\mp\frac12$. Theorem~\ref{thm:T} gives $3p-(n+q+h)=3\theta-80$ up to the cap
$g=6N$, which it reaches at $\theta=\frac{86}3$. From $\theta=29$ the hypotheses of Theorem~\ref{thm:U}
hold. They read $3p\ge L$, $3p\le H$, $2p\le n+q$ and $2h\le p+1$. The bound is
$\min(n,2p+2g+1)-g-2p-1$, which is $6N$ for $\theta\le30$ and $66N-2p-1$ above. In each case the value of the
theorem is the value in the table minus at most one unit, from the integer rounding. The same kind of check,
with linear inequalities in $N$, $p$ and the parameters, gives the bounds of Section~\ref{ssec:families} on
their intervals, and Appendix~\ref{app:windows} gives the others. This proves Theorem~\ref{thm:profile-D}.
The loss of $2$ in its statement covers the integer rounding of all the bounds.